\MakeAbstract{
    建筑幕墙是现代建筑的重要外围护结构，其服役状态直接影响建筑安全、耐久性和运维质量。现实工程中，幕墙巡检长期依赖人工经验，存在效率低、主观性强、难以处理海量图像和难以形成系统化报告等问题。已有计算机视觉方法能够完成裂缝、锈蚀、破损等局部缺陷检测，但多数方案停留在单模型、单图像、单缺陷识别层面，难以解决“从项目级图像资产到建筑级韧性报告”的完整业务链路问题。本文面向这一现实痛点，设计并实现一套基于 AI Agent 的建筑幕墙韧性评估软件系统。

    本文工作的重点有别于重新提出某个检测模型，而在于构建一个能够把既有视觉检测能力组织起来、让 Agent 参与任务规划与信息整合的工程框架。系统首先管理项目基础信息和幕墙图像资产，然后由分类 Agent 根据图像内容自主判断需要执行的原子检测能力，包括金属锈蚀、石材裂缝、石材污渍、玻璃平整度和玻璃爆裂。检测完成后，系统进一步提出分层信息整合策略：先将每张图像的多源检测结果生成单图总结，再将所有单图总结、项目背景、图像分组和人工复核意见融合为项目级韧性评估报告。该策略将大语言模型从简单文本生成工具扩展为“路由、压缩、汇总”的智能节点，缓解了海量检测数据直接输入模型带来的上下文过载和重点淹没问题。

    在工程实现上，系统采用 React 前端、Golang 后端、gRPC 服务通信、RocketMQ 异步任务触发、WebSocket 实时进度推送、MySQL 结构化存储、Redis 缓存和 MinIO 对象存储。工程链路为 Agent 能力提供稳定运行环境，使分类、检测、总结和报告生成过程可追踪、可回调、可恢复和可展示。最终系统形成从图像上传、智能检测、人工复核到报告生成的完整闭环，为建筑幕墙运维提供了一种兼具工程可落地性和智能决策能力的解决方案。
}{AI Agent，建筑幕墙，韧性评估，任务编排}

\MakeAbstractEng{
    Building curtain walls are important envelope structures in modern buildings, and their service condition directly affects building safety, durability, and operation quality. In real engineering practice, curtain wall inspection has long relied on manual experience, leading to low efficiency, strong subjectivity, difficulty in processing massive image data, and difficulty in producing systematic reports. Existing computer vision methods can detect local defects such as cracks, corrosion, and breakage, but most solutions remain at the level of single-model, single-image, and single-defect recognition. They are still insufficient for solving the complete business workflow from project-level image assets to building-level resilience assessment reports. To address this practical pain point, this thesis designs and implements an AI Agent-based software system for building curtain wall resilience assessment.

    The focus of this work is not to propose a new detection model, but to build an engineering framework that organizes existing visual detection capabilities and enables Agents to participate in task planning and information integration. The system first manages project profile information and curtain wall image assets. Then, a classification Agent autonomously determines which atomic detection capabilities should be executed according to image content, including metal corrosion detection, stone crack detection, stone stain detection, glass flatness detection, and glass spalling detection. After detection is completed, the system further proposes a hierarchical information integration strategy: multi-source detection results of each image are first summarized into an image-level summary, and then all image-level summaries, project background, image groups, and human review comments are integrated into a project-level resilience assessment report. This strategy extends large language models from simple text generation tools into intelligent nodes for routing, compression, and aggregation, alleviating the context overload and key-information dilution caused by directly feeding massive detection data into the model.

    In terms of engineering implementation, the system adopts a React frontend, Golang backend services, gRPC service communication, RocketMQ-based asynchronous task triggering, WebSocket-based real-time progress pushing, MySQL structured storage, Redis caching, and MinIO object storage. The engineering workflow provides a stable runtime environment for Agent capabilities, making classification, detection, summarization, and report generation traceable, callback-driven, recoverable, and displayable. The final system forms a complete closed loop from image upload, intelligent detection, and human review to report generation, providing a solution for building curtain wall operation and maintenance that combines engineering feasibility with intelligent decision-making capability.
}{AI Agent, Building Curtain Wall, Resilience Assessment, Task Orchestration}
